\section{Πρόβλημα 5}

Θεωρούμε παίγνιο $n$ παικτών όπου ο $i$-οστός παίκτης έχει σύνολο στρατηγικών
$S_i = \{s_{i1}, s_{i2}, \ldots, s_{im_i}\}$. Θα αποδείξουμε ή θα ανταποδείξουμε
τα ακόλουθα.

\subsection{Ερώτηση (i)}

Αν το παίγνιο έχει αμιγείς Nash ισορροπίες τότε οποιαδήποτε κατανομή επάνω στα
προφίλ στρατηγικών που αντιστοιχούν σε αμιγείς Nash ισορροπίες είναι ένα
(Coarse) Correlated Equilibrium.

Έστω $\sigma$ μια αυθαίρετη κατανομή στα προφίλ στρατηγικών που αποτελούν
αμιγείς ισορροπίες. Θέλουμε να δείξουμε ότι
$$
    \forall i \in \mathbb{N} \colon \forall s_i, s_i' \in S_i \colon
        \mathbb{E}_{s \sim \sigma} [c_i(s_i, s_{-i}) | s_i] \leq
        \mathbb{E}_{s \sim \sigma} [c_i(s_i', s_{-i}) | s_i] \,.
$$

Δεδομένου ότι κάθε $s$ αντιστοιχεί σε αμιγή ισορροπία, ισχύει ότι
$c_i(s_i, s_{-i}) \leq c_i(s_i', s_{-i})$ και άρα
$$
    \mathbb{E}_{s \sim \sigma} [c_i(s_i, s_{-i}) | s_i] \leq
    \mathbb{E}_{s \sim \sigma} [c_i(s_i', s_{-i}) | s_i] \,.
$$

\subsection{Ερώτηση (ii)}

Αν το παίγνιο έχει ένα Correlated Equilibrium $Eq$ όπου για κάθε $i, j$ το $s_{ij}$
εμφανίζεται το πολύ σε ένα από τα προφίλ στρατηγικών με θετική πιθανότητα στο $Eq$,
τότε κάθε προφίλ στρατηγικών με θετική πιθανότητα στο $Eq$ είναι αμιγής Nash
ισορροπία.

Έστω $s$ ένα προφίλ με θετική πιθανότητα στο $Eq$ και $s_i$ η στρατηγική του
παίκτη $i$ στο $s$. Αφού το $s_i$ εμφανίζεται σε ακριβώς ένα προφίλ με θετική
πιθανότητα (το $s$), η δεσμευμένη κατανομή $\sigma(\cdot \mid s_i)$ είναι
συγκεντρωμένη στο $s$, άρα
$\mathbb{E}_{s \sim \sigma}[c_i(s_i, s_{-i}) \mid s_i] = c_i(s_i, s_{-i})$
και
$\mathbb{E}_{s \sim \sigma}[c_i(s_i', s_{-i}) \mid s_i] = c_i(s_i', s_{-i})$.
Επομένως, ο ορισμός του Correlated Equilibrium δίνει για κάθε $s_i' \in S_i$:
$$
    c_i(s_i, s_{-i}) \leq c_i(s_i', s_{-i}) \,,
$$
που είναι ακριβώς η συνθήκη αμιγούς Nash ισορροπίας για το $s$.

\subsection{Ερώτηση (iii)}

Αν το παίγνιο έχει Coarse Correlated Equilibrium $Eq$ όπου για κάθε $i, j$ το
$s_{ij}$ εμφανίζεται το πολύ σε ένα από τα προφίλ στρατηγικών με θετική πιθανότητα
στο $Eq$, τότε κάποιο από τα προφίλ στρατηγικών με θετική πιθανότητα στο $Eq$
είναι αμιγής Nash ισορροπία. Θα ανταποδείξουμε αυτό.

Θεωρούμε παίγνιο 2 παικτών όπου ο παίκτης 1 έχει στρατηγικές $\{A, B, C\}$ και
ο παίκτης 2 έχει στρατηγικές $\{D, E, F\}$

\begin{center}
\begin{tabular}{cccc}
        & $D$ & $E$ & $F$ \\
  $A$   & (3, 3)   & (5, 3)   & (1, 5)   \\
  $B$   & (1, 3)   & (3, 3)   & (5, 1)   \\
  $C$   & (3, 3)   & (3, 3)   & (3, 3)   \\
\end{tabular}
\,.
\end{center}

Θεωρούμε την ομοιόμορφη κατανομή $\sigma$ πάνω στα προφίλ
$\{(A,D),\, (B,E),\, (C,F)\}$. Κάθε στρατηγική εμφανίζεται ακριβώς σε ένα από τα
τρία προφίλ.

Ο αναμενόμενος κόστος κάθε παίκτη είναι
$\mathbb{E}_\sigma[c_i] = \frac{3+3+3}{3} = 3$.
Για κάθε εναλλακτική στρατηγική $s'$ του παίκτη 1:
$$
    \mathbb{E}_\sigma[c_1(A, s_2)] = \tfrac{3+5+1}{3} = 3, \quad
    \mathbb{E}_\sigma[c_1(B, s_2)] = \tfrac{1+3+5}{3} = 3, \quad
    \mathbb{E}_\sigma[c_1(C, s_2)] = \tfrac{3+3+3}{3} = 3 \,.
$$
Συμμετρικά, για τον παίκτη 2:
$$
    \mathbb{E}_\sigma[c_2(s_1, D)] = \tfrac{3+3+3}{3} = 3, \quad
    \mathbb{E}_\sigma[c_2(s_1, E)] = \tfrac{3+3+3}{3} = 3, \quad
    \mathbb{E}_\sigma[c_2(s_1, F)] = \tfrac{5+1+3}{3} = 3 \,.
$$
Άρα κανένας παίκτης δεν ωφελείται μεταβαίνοντας σε σταθερή στρατηγική, οπότε
το $\sigma$ είναι CCE.

Ωστόσο, κανένα από τα τρία προφίλ δεν είναι αμιγής Nash ισορροπία:
\begin{itemize}
    \item Στο $(A,D)$: ο παίκτης 1 προτιμά $B$, αφού $c_1(B,D) = 1 < 3 = c_1(A,D)$.
    \item Στο $(B,E)$: ο παίκτης 2 προτιμά $F$, αφού $c_2(B,F) = 1 < 3 = c_2(B,E)$.
    \item Στο $(C,F)$: ο παίκτης 1 προτιμά $A$, αφού $c_1(A,F) = 1 < 3 = c_1(C,F)$.
\end{itemize}
